%\usepackage[mtbold]{mathtime}
\usepackage{enumerate}
\usepackage{xypic}

\setlength{\oddsidemargin}{0in}
\setlength{\evensidemargin}{0in}
\setlength{\textwidth}{6.5in}
\setlength{\topmargin}{0in}
\setlength{\textheight}{8.5in}

\setcounter{topnumber}{3}%
\def\topfraction{.7}%
\setcounter{bottomnumber}{1}
\def\bottomfraction{.3}
\setcounter{totalnumber}{5}%
\def\textfraction{.1}% was .2
\def\floatpagefraction{.7}% was .7
\setcounter{dbltopnumber}{2}
\def\dbltopfraction{.7}
\def\dblfloatpagefraction{.5}
%\addtolength{\parskip}{\baselineskip}

\long\def\comment#1{\ignorespaces}

\pagestyle{myheadings}

%% Parameters to the \lecture command are the lecture number, the
%% date, and the scribe's name.
\newcommand{\lecture}[3]{\newpage\markboth{\textit{Lecture #1, #2}}{\textit{Lecture #1, #2}}
{
\begin{flushleft}
\LARGE{\bf Lecture #1}\\
\LARGE{\bf #2}\\
\end{flushleft}}
\bigskip}

\newenvironment{example}{\begin{quote}}{\end{quote}}
\newenvironment{letter}[1]{\begin{enumerate}[a.]\item[#1.]}{\end{enumerate}}
\newenvironment{letters}{\begin{enumerate}[a.]}{\end{enumerate}}
\newcommand{\xmp}[1]{\textit{#1}}
\newcommand{\good}{\emph{Good:}}
\newcommand{\bad}{\emph{Bad:}}
\newcommand{\ok}{\emph{OK:}}
\newcommand{\better}{\emph{Better:}}
\newcommand{\wrong}{\emph{Wrong:}}
\newcommand{\correct}{\emph{Right:}}

\newtheorem{theorem}{Theorem}
\newtheorem{lemma}[theorem]{Lemma}
\newtheorem{corollary}[theorem]{Corollary}

% Make captions with bold figure/table name and number.
\makeatletter
\long\def\@makecaption#1#2{%
  \vskip\abovecaptionskip
  \sbox\@tempboxa{\small\textbf{#1:} #2}%
  \ifdim \wd\@tempboxa >\hsize
    {\small\textbf{#1:} #2}\par
  \else
    \global \@minipagefalse
    \hb@xt@\hsize{\hfil\box\@tempboxa\hfil}%
  \fi
  \vskip\belowcaptionskip}
\makeatother

\newcommand{\figcaption}[1]{\caption[]{#1}}
\newcommand{\figpart}[1]{\textbf{(#1)}~\ignorespaces}

\newcommand{\subheading}[1]{\subsubsection*{#1}}
\newcommand{\subsubheading}[1]{\paragraph*{#1.}}

\newcommand{\defn}[1]{\emph{#1}}

\newcommand{\eqnref}[1]{(\ref{#1})}
\renewcommand{\cases}[1]{\left\{ \begin{array}{ll}#1\end{array}\right.}
\newcommand{\cif}[1]{\mbox{if $#1$}}

\newcommand{\prob}[1]{\Pr\left\{ #1 \right\}}
\renewcommand{\hat}{\widehat}
\newcommand{\elems}[3]{#1, \lb #2, \lb \ldots, #3}
\renewcommand{\vector}[1]{(#1)}
\newcommand{\evector}[3]{\vector{\elems{#1}{#2}{#3}}}
\newcommand{\seq}[1]{\langle#1\rangle}
\newcommand{\set}[1]{\{#1\}}
\newcommand{\eset}[3]{\set{\elems{#1}{#2}{#3}}}
\def\ceil#1{\left\lceil#1\right\rceil}
\def\floor#1{\left\lfloor#1\right\rfloor}
\newcommand{\abs}[1]{\left|#1\right|}
\newcommand{\card}[1]{\left|#1\right|}
\newcommand{\smallcard}[1]{|#1|}
\newcommand{\proof}{\noindent\textit{Proof:}\hspace*{1em}}
\newcommand{\proofof}[1]{\noindent\textit{Proof of #1:}\hspace*{1em}}
\newcommand{\proofsketch}{\noindent\textit{Proof sketch:}\hspace*{1em}}
\newcommand{\proofnote}[1]{\noindent\textit{Proof:}\hspace*{.1em}\footnote{#1}\hspace*{0.5em}}
\def\twodots{\mathinner{\ldotp\ldotp}}
\newcommand{\range}{\twodots}
\newcommand{\inv}{^{-1}}
\newcommand{\transpose}{^{\rm T}}
\newcommand{\ptranspose}{^{\prime{\rm T}}}
\newcommand{\oneone}{\stackrel{\mbox{\scriptsize 1-1}}{\rightarrow}}
\newcommand{\xor}{\oplus}
\newcommand{\bigxor}{\bigoplus}
\def\degrees{^{\circ}}

\DeclareMathSymbol{\umod}{\mathbin}{operators}{"23}

\def\qedbox{\hspace{-0.3em}\mbox{\vrule height 5pt width 5pt}}
\def\qed{\unskip\parfillskip=0pt~~~~\hfill\llap{\vrule height 5pt width 5pt}\vskip11pt\parfillskip=0pt plus 1fil\relax}
\def\qednoproof{\unskip\parfillskip=0pt~~~~\hfill\llap{\vrule height 5pt width 5pt}}

\newenvironment{eqnarrayreason}{\[\begin{array}{rcll}}{\end{array}\]}
\newcommand{\eqnline}{\\[.75ex]}
\newcommand{\bigeqnline}{\\[2.5ex]}

\newlength{\savearraycolsep}
\newcommand{\narrowarray}[1]{\setlength{\savearraycolsep}{\arraycolsep}%
\setlength{\arraycolsep}{#1\arraycolsep}}
\newcommand{\normalarray}{\setlength{\arraycolsep}{\savearraycolsep}}
\newlength{\savetabcolsep}
\newcommand{\narrowtabular}[1]{\setlength{\savetabcolsep}{\tabcolsep}%
\setlength{\tabcolsep}{#1\tabcolsep}}
\newcommand{\normaltabular}{\setlength{\tabcolsep}{\savetabcolsep}}

\newcommand{\noarraystretch}{\def\arraystretch{}}
\newcommand{\matx}[2]{{\noarraystretch\left[\begin{array}{*{#1}{c}}#2\end{array}\right]}}
\newcommand{\block}[2]{{\noarraystretch\count200=#1%
\advance\count200 by -1\left[\begin{array}{*{\count200}{c|}c}#2\end{array}\right]}}
\newcommand{\hatstrut}{\rule{0pt}{3ex}}

\newcommand{\bstrut}{\rule{0pt}{2.4ex}}
\newcommand{\bline}{\hline\bstrut}
\newlength{\blockdimwidth}
\newcommand{\bm}[1]{\makebox[\blockdimwidth]{$#1$}}
\newcommand{\bt}[1]{\makebox[\blockdimwidth]{\footnotesize$#1$}}
\newcommand{\bs}[1]{\mbox{\bstrut\footnotesize$#1$}}
\newcommand{\blockdims}[6]{{\noarraystretch\setlength{\blockdimwidth}{#1}%
\begin{array}{cc@{\hspace{0em}}c} &%
\begin{array}{c*{#2}{c}}#4\end{array} & \\%
#3 &
\count200=#2\advance\count200 by -1%
\left[\begin{array}{*{\count200}{c|}c}#5\end{array}\right] &%
\begin{array}{l}#6\end{array}%
\end{array}}}
\newcommand{\blockdimscol}[6]{{\noarraystretch\setlength{\blockdimwidth}{#1}%
\begin{array}{cc@{\hspace{0em}}c} &%
\begin{array}{cc}#4\end{array} & \\%
#3 &
\left[\begin{array}{c}#5\end{array}\right] &%
\begin{array}{l}#6\end{array}%
\end{array}}}
\newcommand{\blockdimstoponly}[5]{{\noarraystretch\setlength{\blockdimwidth}{#1}%
\begin{array}{cc} &%
\begin{array}{c*{#2}{c}}#4\end{array}\\%
#3 &
\count200=#2\advance\count200 by -1%
\left[\begin{array}{*{\count200}{c|}c}#5\end{array}\right]%
\end{array}}}
\newcommand{\blockdimsmatxonly}[5]{{\noarraystretch\setlength{\blockdimwidth}{#1}%
\begin{array}{c@{\hspace{0em}}c}%
\begin{array}{c*{#2}{c}}#3\end{array} & \\%
\count200=#2\advance\count200 by -1%
\left[\begin{array}{*{\count200}{c|}c}#4\end{array}\right] &%
\begin{array}{l}#5\end{array}%
\end{array}}}
\newcommand{\blockdimstopmatxonly}[4]{{\noarraystretch\setlength{\blockdimwidth}{#1}%
\begin{array}{c@{\hspace{0em}}c}%
\begin{array}{c*{#2}{c}}#3\end{array} & \\%
\count200=#2\advance\count200 by -1%
\left[\begin{array}{*{\count200}{c|}c}#4\end{array}\right]%
\end{array}}}

\usepackage{url}
\urlstyle{rm}
